\subsubsection{A Universal Kaluza-Klein Graviton Cascade Rate --- arXiv:2609.36234}

\textbf{Authors:} Vincent S. H. Lee, Lisa Randall, Marcos Riojas

\paragraph{主な主張}
5D KK graviton の三点相互作用では massless 5D graviton の運動方程式に由来して $m_n^2$ に比例する leading term が相殺され、near-threshold では振幅が $q^2$ に比例するため、二体位相空間まで含めた cascade rate は $\Gamma\propto q^5$ となる；これは従来の Dark Dimension phenomenology で用いられた見積もりに対して $(q/m_n)^4$ の追加抑制を与える \cite{Lee:2026KKCascade}。

\paragraph{新規性}
一般の warp factor、brane と scalar の backreaction、さらに non-minimal な二階微分 scalar--gravity coupling を含む広い 5D 理論でこの cancellation を示し、scalar-induced KK-number violation を warp factor に吸収することで Randall--Sundrum 模型で知られていた $q^5$ scaling の普遍性を示した。

\paragraph{位置付け}
rapid KK cascade を前提とする Dark Dimension graviton dark matter の中心的仮定を直接再検討し、既存模型では Standard Model への decay constraint を回避できないことを示す結果である；高次元化では、各 decay channel の $q^5$ suppression が残るか、増加する KK state multiplicity が inclusive rate をどこまで補うかを調べる際の基準点になる。

\paragraph{Introduction の流れ}
KK number が破れれば重い KK mode が軽い mode へ崩壊できることは古くから知られ、bulk scalar によって extra-dimensional momentum conservation を破る具体例も提案されていた \cite{Mohapatra:2003KKDecay}；Dienes--Thomas は多数の不安定成分からなる Dynamical Dark Matter の枠組みを構成し、KK tower をその自然な実現例として位置付けた \cite{Dienes:2011DDM1,Dienes:2011DDM2}。その後 GMOV はこの考えを Dark Dimension の graviton KK tower に適用し \cite{Gonzalo:2022DarkGravitons,Vafa:2024DarkSector}、宇宙論・天体物理制限を回避するには十分速い cascade が必要であることが具体的に議論された \cite{Obied:2023DecayingGravitons,LawSmith:2023AstroGravitons}。一方で decay amplitude の運動量スケールを $M=m_n$ と置く見積もりと $M=q$ とする Randall--Sundrum 計算 \cite{deGiorgi:2021Spin2} が食い違っており、本論文は全ての cubic gravitational terms を保つと leading term が相殺され、普遍的に $M=q$ となることを示してこの曖昧さを解消する \cite{Lee:2026KKCascade}。
